
\subsection{Grundlegende Definitionen}

\begin{definition}
  Sei $M$ eine Menge. 
  \begin{enumerate}[(a)]
  \item Eine Abbildung $\_\cdot\_:M\times M \to M$ nennen wir auch \begriff{Verknüpfung}.
  \item Ein \begriff{Magma} ist eine Menge $M$ zusammen mit einer Verknüpfung. 
  \item Eine Verknüpfung $\_\cdot\_:M\times M\to M$ heißt \begriff{assoziativ}, wenn gilt:
    \begin{align*}
      \forall x,y,z\in M. (x\cdot y)\cdot z = x\cdot (y\cdot z)
    \end{align*}
  \item Ein Magma mit assoziativer Verknüpfung heißt \begriff{Halbgruppe}
  \item Sei $M$ ein Magma mit Verknüpfung $\_\cdot\_$. Ein Element $e\in M$ heißt \begriff{linksneutral}, wenn
    \begin{align*}
      \forall x\in M. e\cdot x=x
    \end{align*}
    und \begriff{rechtsneutral}, wenn 
    \begin{align*}
      \forall x\in M. x\cdot e=x
    \end{align*}
    gilt. $e$ heißt \begriff{neutrales Element}, wenn $e$ links- und rechtsneutral ist.
  \item Ein Magma $(M,\_\cdot\_)$ zusammen mit einem neutralen Element $e\in M$ heißt \begriff{Monoid}.
  \item Sei $M$ ein Monoid und $x,y\in M$. $y$ heißt \begriff{rechtsinvers} zu $x$, wenn $x\cdot y=e$ gilt. $y$ heißt \begriff{linksinvers} zu $x$, wenn $y\cdot x=e$ gilt. Wenn $y$ link- und rechtsinvers zu $x$ ist, sagen wir $y$ ist \begriff{invers} zu $x$.
  \item Ein Monoid $(M,\_\cdot\_,e)$ heißt \begriff{Gruppe}, wenn es zu jedem Element ein inverses gibt.
  \end{enumerate}
\end{definition}

\begin{bemerkung}
  \begin{enumerate}[(a)]
  \item 
  \end{enumerate}
\end{bemerkung}
